\documentclass[11pt]{article}
\usepackage[margin=1in]{geometry}
\usepackage{amsmath,amssymb,amsthm}
\usepackage{graphicx}
\usepackage{hyperref}

\newtheorem{theorem}{Theorem}
\newtheorem{lemma}{Lemma}
\newtheorem{corollary}{Corollary}
\newcommand{\R}{\mathbb R}

\title{D.c. extensions from finite-codimensional polyhedral cylinders}
\author{Partial answer to Remark 2.4(a) of arXiv:0810.1433}
\date{12 August 2026}

\begin{document}
\maketitle

\noindent\textbf{Status.} \texttt{candidate\_partial\_likely\_valid}.

\begin{abstract}
Vesel\'y and Zaj\'{\i}\v{c}ek ask whether a power-type-$2$ uniformly convex
renorming assumption can be removed or weakened in their d.c.-extension
theorem. We remove all renorming assumptions for finite-codimensional
polyhedral cylinders. If $T:X\to E$ is onto a finite-dimensional space,
$K\subset E$ is a full-dimensional closed convex polyhedron, and
$C=T^{-1}(K)$, then every d.c. mapping $F:C\to Y$ for which $F$ and
one control function are Lipschitz on bounded parts of $C$ has a d.c.
extension to $X$. The proof replaces the source's radial retraction by a
finite-rank piecewise-affine retraction. The ambient normed space $X$ and
the target $Y$ are otherwise arbitrary.
\end{abstract}

\section{Source problem}

For a convex set $D$ in a normed space, a mapping $G:D\to Y$ is d.c.
when there is a continuous convex $g:D\to\R$ such that
$y^*\circ G+g$ is convex for every $y^*\in Y^*$ with
$\|y^*\|\le1$. Such a $g$ is a \emph{control function}.

Theorem 2.3 of Vesel\'y--Zaj\'{\i}\v{c}ek \cite{source} gives equivalent
extension criteria under the assumption that $X$ admits an equivalent norm
with modulus of convexity of power type $2$. Remark 2.4(a), reproduced in
Figure~\ref{fig:source}, asks whether that assumption can be omitted or
essentially weakened.

The source uses this geometry to prove that the radial retraction onto
$\overline C$ is d.c.; its remaining composition argument is geometric-free.
The following result supplies a different retraction.

\section{Idea of the proof}

A finite-dimensional polyhedron has a nearest-point projection which is
piecewise affine on finitely many polyhedral regions. Finite piecewise-affine
maps are d.c. with global Lipschitz controls. If the polyhedron is observed
through a finite-dimensional quotient $T:X\to E$, a linear right inverse
lifts that projection to a d.c. retraction of the arbitrary normed space
$X$ onto $T^{-1}(K)$. The retraction sends bounded sets to bounded sets,
so the source's local composition lemma applies ball by ball.

\section{The piecewise-affine retraction}

We use the following standard finite-dimensional fact in a form that records
the control needed below.

\begin{lemma}[Finite piecewise-affine maps]\label{lem:pwa}
Let $E$ and $Z$ be finite-dimensional normed spaces. Every continuous map
$H:E\to Z$ which is affine on each cell of a finite polyhedral subdivision
is globally Lipschitz and d.c. with a globally Lipschitz control function.
\end{lemma}

\begin{proof}
It suffices to treat a scalar component $h$. Let
$\lambda_1=0,\ldots,\lambda_N=0$ be the finite collection of affine
hyperplanes containing the facets of the subdivision. Across a common facet,
continuity forces the jump of the affine gradient of $h$ to be a scalar
multiple of the normal gradient of the corresponding $\lambda_i$. For a
sufficiently large $M$,
\[
 g=M\sum_{i=1}^N |\lambda_i|
\]
has a normal gradient jump which dominates, with either sign, every such jump
of $h$. The standard facet criterion for convex piecewise-affine functions
then shows that $g+h$ and $g-h$ are convex. Thus $g$ controls $h$, and
both $h$ and $g$ are globally Lipschitz.

Choose coordinates on $Z$, apply the scalar argument to every component,
and take a sufficiently large positive sum of the scalar controls. Writing a
coefficient $a$ with $|a|\le c$ as a positive combination of the controls
for $h$ and $-h$ shows that this single sum controls every member of the
dual unit ball of $Z$.
\end{proof}

\begin{lemma}[Polyhedral-cylinder retraction]\label{lem:retraction}
Let $X$ be a normed space, let $E$ be finite-dimensional, and let
$T:X\to E$ be a continuous linear surjection. If $K\subset E$ is a
nonempty closed convex polyhedron, then $C=T^{-1}(K)$ admits a globally
Lipschitz d.c. retraction $R:X\to C$ with a globally Lipschitz control.
In particular, $R$ maps bounded sets to bounded sets.
\end{lemma}

\begin{proof}
Equip $E$ with an auxiliary Euclidean norm and let
$\pi_K:E\to K$ be nearest-point projection. It is $1$-Lipschitz. For
each face $G$ of $K$, the points whose projection lies in the relative
interior of $G$ form a polyhedral region, and on that region $\pi_K$ is
orthogonal projection onto the affine hull of $G$. Since $K$ has finitely
many faces, $\pi_K$ is a finite continuous piecewise-affine map. It is d.c.
with a Lipschitz control by Lemma~\ref{lem:pwa}.

Because $E$ is finite-dimensional, $T$ has a continuous linear right
inverse $S:E\to X$. Define
\[
 R(x)=x+S\bigl(\pi_K(Tx)-Tx\bigr).                           \tag{1}\label{eq:R}
\]
Linear pre- and postcomposition and addition of an affine map preserve d.c.
mappings and Lipschitz controls, so $R$ has the asserted regularity. Also
\[
 TR(x)=\pi_K(Tx)\in K,
\]
while $R(x)=x$ whenever $Tx\in K$. Thus $R$ is a retraction onto $C$.
Its global Lipschitz bound implies bounded-set preservation.
\end{proof}

\section{Extension theorem without ambient renorming}

For the next statement, “Lipschitz on bounded parts of $C$” means Lipschitz
on $C\cap rB_X$ for every $r>0$.

\begin{theorem}[Polyhedral extension]\label{thm:main}
Let $X,Y$ be arbitrary normed spaces, $E$ finite-dimensional,
$T:X\to E$ a continuous linear surjection, and $K\subset E$ a
full-dimensional closed convex polyhedron. Put $C=T^{-1}(K)$. Suppose
$F:C\to Y$ is d.c. and has a control function $f:C\to\R$ such that both
$F$ and $f$ are Lipschitz on bounded parts of $C$. Then $F$ admits a
d.c. extension to all of $X$. No renorming hypothesis on $X$ is required.
\end{theorem}

\begin{proof}
Let $R:X\to C$ be the retraction from Lemma~\ref{lem:retraction} and set
\[
 \widehat F=F\circ R.
\]
This extends $F$. To prove it is d.c., use Lemma 1.12 of the source. Let
$A_n=n\,\operatorname{int}B_X$. Since $R$ is globally Lipschitz, choose
$M_n<\infty$ such that
\[
 R(A_n)\subset B_n:=C\cap M_nB_X.
\]
Each $B_n$ is convex. By hypothesis, $F|_{B_n}$ is Lipschitz and d.c. with
the Lipschitz control $f|_{B_n}$. The map $R$ is d.c., the open convex
sets $A_n$ increase to $X$, and $R(A_n)\subset B_n$. Every hypothesis
of the source's composition lemma is therefore satisfied, so $F\circ R$ is
d.c. on $X$.
\end{proof}

\begin{corollary}
The conclusion holds, on every normed space, when $C$ is a halfspace, a
finite-codimensional slab or box, or more generally the inverse image of any
finite-dimensional closed convex polyhedron under a quotient map.
\end{corollary}

\section{Scope and obstruction to a full upgrade}

Theorem~\ref{thm:main} gives a genuine class of infinite-dimensional spaces
and domains for which the source's renorming assumption disappears. It does
strengthen the local hypothesis by requiring $F$ and one control to be
Lipschitz on bounded parts of $C$. It does not settle Theorem 2.3 for an
arbitrary convex $C$.

The remaining difficulty is not the final composition step but the existence
of a suitable d.c. retraction. For a general $C$, the source's radial map
multiplies a scalar d.c. function by the full infinite-dimensional vector
$x$. Flat directions of arbitrary norms prevent the obvious radial convex
controls from dominating the mixed bilinear curvature. Eight focused routes,
including scalar splitting, UMD geometry, and counterexample reductions, are
recorded in the accompanying attempt log.

Bounded arXiv searches through 12 August 2026 found no later resolution of
Remark 2.4(a) and no statement of Theorem~\ref{thm:main}. Novelty confidence
is moderate-low because the retraction argument is elementary and may be
unindexed folklore. Expert review should focus on Lemma~\ref{lem:pwa} and the
scope comparison with the source theorem.

\begin{thebibliography}{9}
\bibitem{source}
L. Vesel\'y and L. Zaj\'{\i}\v{c}ek,
\emph{On extensions of d.c. functions and convex functions},
arXiv:0810.1433; Journal of Convex Analysis 17 (2010), 427--440.
\end{thebibliography}

\clearpage
\begin{figure}[p]
\centering
\includegraphics[width=\textwidth]{figures/open_problem_crop.png}
\caption{Source page 7: the radial-retraction step and Remark 2.4(a).}
\label{fig:source}
\end{figure}

\end{document}
